\documentclass[11pt,a4paper]{article}

\usepackage{fontspec}
\usepackage[croatian]{babel}
\usepackage[a4paper,margin=2.2cm,top=2.4cm,bottom=2.4cm]{geometry}
\usepackage{xcolor}
\usepackage{array,tabularx,booktabs,colortbl}
\usepackage{enumitem}
\usepackage{fancyvrb}
\usepackage[most]{tcolorbox}
\usepackage{titlesec}
\usepackage{fancyhdr}
\usepackage{lastpage}
\usepackage{soul}
\usepackage{needspace}
\usepackage[hidelinks]{hyperref}

\setmainfont{IBMPlexSans}[
  Extension=.otf, UprightFont=*-Regular, ItalicFont=*-Italic,
  BoldFont=*-SemiBold, BoldItalicFont=*-SemiBoldItalic]
\setsansfont{IBMPlexSans}[
  Extension=.otf, UprightFont=*-Regular, ItalicFont=*-Italic,
  BoldFont=*-SemiBold, BoldItalicFont=*-SemiBoldItalic]
\setmonofont{IBMPlexMono}[
  Extension=.otf, UprightFont=*-Regular, ItalicFont=*-Italic,
  BoldFont=*-SemiBold, BoldItalicFont=*-SemiBoldItalic, Scale=0.9]

\definecolor{ink}{HTML}{0C4685}
\definecolor{ink2}{HTML}{6E7BAA}
\definecolor{tint}{HTML}{EAEEF6}
\definecolor{tekst}{HTML}{1A222E}
\definecolor{muted}{HTML}{6B7484}
\definecolor{pozadina}{HTML}{F6F7F9}
\definecolor{linije}{HTML}{D9DEE7}

\color{tekst}
\setlength{\parindent}{0pt}
\setlength{\parskip}{0.55em}
\linespread{1.12}
\hypersetup{colorlinks=true,urlcolor=ink,linkcolor=ink}
\sethlcolor{tint}

\titleformat{\section}{\Large\bfseries\color{ink}}{}{0pt}{}
\titleformat{\subsection}{\large\bfseries\color{ink}}{}{0pt}{}
\titlespacing*{\section}{0pt}{1.4em}{0.4em}
\titlespacing*{\subsection}{0pt}{1.1em}{0.3em}

\pagestyle{fancy}
\fancyhf{}
\renewcommand{\headrulewidth}{0pt}
\fancyhead[L]{\footnotesize\color{muted}Programiranje za web - priprema razvojnog okruženja}
\fancyfoot[R]{\footnotesize\color{muted}Stranica \thepage\ od \pageref*{LastPage}}

\newcommand{\code}[1]{{\ttfamily\small #1}}
\newcolumntype{L}{>{\raggedright\arraybackslash}X}
\newcolumntype{P}[1]{>{\raggedright\arraybackslash}p{#1}}
\renewcommand{\arraystretch}{1.35}

\newenvironment{kod}
  {\VerbatimEnvironment
   \begin{tcolorbox}[colback=pozadina,colframe=linije,boxrule=0.4pt,arc=2pt,
     left=6pt,right=6pt,top=3pt,bottom=3pt]
   \begin{Verbatim}[fontsize=\small]}
  {\end{Verbatim}\end{tcolorbox}}

\newenvironment{ispis}
  {\VerbatimEnvironment
   \begin{tcolorbox}[colback=pozadina,colframe=linije,boxrule=0.4pt,arc=2pt,
     left=6pt,right=6pt,top=3pt,bottom=3pt]
   \begin{Verbatim}[fontsize=\footnotesize]}
  {\end{Verbatim}\end{tcolorbox}}

\newcommand{\zaglavlje}{\rowcolor{tint}}

\begin{document}

{\raggedright\hyphenpenalty=10000\color{ink}\fontsize{22}{26}\selectfont\bfseries Programiranje za web - priprema razvojnog okruženja\par}
\vspace{0.3em}
{\color{muted}28. rujna 2026. \textperiodcentered{} Josip Torić\par}
\vspace{0.4em}

\section{Uvod}

Do prvog predavanja \textbf{8.\ listopada 2026.} na svom računalu instalirajte sve alate
s popisa i pokrenite radni primjer iz mape \code{pzw-provjera}. Ako svaka naredba ispiše
ono što je navedeno, okruženje je spremno za cijeli kolegij.

Ako nešto ne radi, pogledajte odjeljak \hyperref[sec:problemi]{\emph{Česti problemi}} na kraju
uputa. Ako to ne pomogne, javite se asistentu Josipu Toriću
(\href{mailto:jtoric@unizd.hr}{jtoric@unizd.hr}) \ul{prije prvog predavanja} kako bismo
osigurali neometan početak nastave.

\section{Popis softvera}

Instalirajte šest alata i ekstenzije iz tablice, svaki prema službenim uputama za svoj
operacijski sustav. Nakon svake instalacije zatvorite i ponovno otvorite terminal, odnosno
cijeli VS Code, da se osvježi PATH.

\begin{center}\small
\begin{tabularx}{\linewidth}{P{2.6cm}P{2.6cm}LP{4.2cm}}
\arrayrulecolor{linije}\toprule
\zaglavlje \textbf{Alat} & \textbf{Verzija} & \textbf{Za što služi} & \textbf{Upute za instalaciju}\\
\midrule
Git & najnovija & verzioniranje koda i predaja projekta & \href{https://git-scm.com/downloads}{git-scm.com/downloads}\\
Python & 3.11 ili novija, preporuka 3.13 & backend u FastAPI-ju & \href{https://www.python.org/downloads/}{python.org/downloads}\\
uv & najnovija & virtualna okruženja i paketi & \href{https://docs.astral.sh/uv/getting-started/installation/}{docs.astral.sh/uv}\\
Docker Desktop & najnovija & PostgreSQL baza u kontejneru, uključuje Docker Compose & \href{https://docs.docker.com/desktop/}{docs.docker.com/desktop}\\
Node.js & 24 LTS, minimalno 22.18 & frontend: Vue, Vite i TypeScript & \href{https://nodejs.org/en/download}{nodejs.org/download}\\
VS Code & najnovija & editor & \href{https://code.visualstudio.com/download}{code.visualstudio.com}\\
Ekstenzije za VS Code & najnovije & Python (Microsoft), Vue - Official, SQLTools + PostgreSQL driver & Extensions panel u VS Code-u, prečac Ctrl+Shift+X\\
\bottomrule
\end{tabularx}
\end{center}

Napomene:
\begin{itemize}[leftmargin=1.2em,itemsep=0.2em,topsep=0pt]
  \item \textbf{Windows:} u Python installeru označite \emph{Add python.exe to PATH}.
        Docker Desktop traži WSL~2. Installer ga obično nudi sam, a ako ne, slijedite
        \href{https://learn.microsoft.com/windows/wsl/install}{upute za WSL}.
  \item \textbf{macOS:} sve osim Docker Desktopa možete instalirati kroz Homebrew:
        \code{brew install git python uv node}.
  \item \textbf{Linux:} umjesto Docker Desktopa dovoljni su
        \href{https://docs.docker.com/engine/install/}{Docker Engine i Compose plugin}.
        Korisnika dodajte u grupu \code{docker}.
\end{itemize}

Nakon instalacije Gita jednom postavite ime i e-mail. Koristite isti e-mail kao na GitHubu:

\begin{kod}
git config --global user.name "Ime Prezime"
git config --global user.email "vas.email@example.com"
\end{kod}

\subsection{Provjera verzija}

U bilo kojem terminalu:

\begin{center}\small
\begin{tabularx}{\linewidth}{P{8cm}L}
\arrayrulecolor{linije}\toprule
\zaglavlje \textbf{Naredba} & \textbf{Očekivani ispis}\\
\midrule
\code{git --version} & \code{git version 2.x}\\
\code{python --version}, na macOS-u i Linuxu \code{python3 --version} & \code{Python 3.11} ili novija\\
\code{uv --version} & \code{uv 0.x.y}\\
\code{docker --version} & \code{Docker version 2x.y.z}\\
\code{docker compose version} & \code{Docker Compose version v2.x} ili novija\\
\code{node --version} & \code{v24.x} ili \code{v22.18} i novija\\
\code{npm --version} & \code{10.x} ili novija\\
\code{code --version} & broj verzije VS Code-a\\
\bottomrule
\end{tabularx}
\end{center}

Brojevi verzija smiju biti i veći od navedenih.

\section{Radni primjer}

Radni primjer se nalazi u mapi \code{pzw-provjera} i sadrži iste slojeve koje ćemo
graditi kroz semestar:

\begin{kod}
pzw-provjera/
  docker-compose.yml    # PostgreSQL baza u Dockeru
  api/                  # FastAPI backend koji čita iz baze
  web/                  # Vue + TypeScript frontend
\end{kod}

Mapu otvorite u VS Code-u (\emph{File → Open Folder}), a naredbe pokrećite u njegovu
terminalu (\emph{Terminal → New Terminal}).

\subsection{1. Baza}

Pokrenite Docker Desktop i pričekajte da u izborniku piše \emph{Engine running}.
U mapi \code{pzw-provjera}:

\begin{kod}
docker compose up -d db
docker compose ps
\end{kod}

\needspace{6\baselineskip}
Prvo pokretanje preuzima \emph{image} od oko 150 MB pa traje minutu-dvije. Očekivani ispis
druge naredbe je red s \code{db} u kojem piše \code{Up} i \code{0.0.0.0:5432->5432/tcp},
na primjer:

\begin{SaveVerbatim}{psispis}
NAME                IMAGE         COMMAND                  SERVICE   CREATED         STATUS         PORTS
pzw-provjera-db-1   postgres:17   "docker-entrypoint.s…"   db        7 seconds ago   Up 6 seconds   0.0.0.0:5432->5432/tcp, [::]:5432->5432/tcp
\end{SaveVerbatim}
\begin{tcolorbox}[colback=pozadina,colframe=linije,boxrule=0.4pt,arc=2pt,
  left=6pt,right=6pt,top=3pt,bottom=3pt]
\resizebox{\linewidth}{!}{\BUseVerbatim{psispis}}
\end{tcolorbox}

\subsection{2. Backend}

\begin{kod}
cd api
uv venv

# aktivacija - Windows (PowerShell)
.venv\Scripts\Activate.ps1
# aktivacija - macOS / Linux
source .venv/bin/activate

uv pip install -r requirements.txt
uvicorn main:app --reload
\end{kod}

Nakon zadnje naredbe očekivani ispis je:

\begin{ispis}
INFO:     Will watch for changes in these directories: ['...\\pzw-provjera\\api']
INFO:     Uvicorn running on http://127.0.0.1:8000 (Press CTRL+C to quit)
INFO:     Started reloader process [9296] using WatchFiles
INFO:     Started server process [16616]
INFO:     Waiting for application startup.
INFO:     Application startup complete.
\end{ispis}

Brojevi procesa u uglatim zagradama razlikovat će se kod vas. U pregledniku otvorite:

\begin{center}\small
\begin{tabularx}{\linewidth}{P{6cm}L}
\arrayrulecolor{linije}\toprule
\zaglavlje \textbf{Adresa} & \textbf{Očekivani ispis}\\
\midrule
\url{http://127.0.0.1:8000/health} & \code{\{"status":"ok","database":"PostgreSQL 17..."\}}\\
\url{http://127.0.0.1:8000/docs} & Swagger UI s endpointom \code{GET /health}\\
\bottomrule
\end{tabularx}
\end{center}

Server ostavite da radi, a za sljedeći korak otvorite novi terminal gumbom \textbf{+}
u panelu terminala.

\subsection{3. Frontend}

Iz mape \code{pzw-provjera}:

\begin{kod}
cd web
npm install
npm run type-check
npm run dev
\end{kod}

\code{npm run type-check} ne smije ispisati nijednu grešku. Nakon \code{npm run dev}
u pregledniku otvorite\\\url{http://localhost:5173}. Očekivani ispis je Vue stranica s naslovom
\textbf{You did it!}

\subsection{Završetak}

Po uspješnoj provjeri okruženja zaustavite oba servera s Ctrl+C u njihovim terminalima
i ugasite bazu naredbom \code{docker compose down} u mapi \code{pzw-provjera}. Na predavanjima ćemo krenuti
graditi projekt iznova.

\needspace{12\baselineskip}
\section{Česti problemi}\label{sec:problemi}

Prije svega zatvorite i ponovno otvorite terminal, odnosno cijeli VS Code. Ako to ne
pomogne, ponovno pokrenite računalo.

\begin{center}\small
\begin{tabularx}{\linewidth}{P{4.6cm}P{3.6cm}L}
\arrayrulecolor{linije}\toprule
\zaglavlje \textbf{Simptom} & \textbf{Uzrok} & \textbf{Rješenje}\\
\midrule
\code{python} nije prepoznat ili otvara Microsoft Store na Windowsima & Python nije u PATH-u &
  Ponovno pokrenite installer i označite \emph{Add python.exe to PATH}; u \emph{Settings →
  Apps → Advanced app settings → App execution aliases} isključite aliase za \code{python.exe}\\
\code{Activate.ps1 cannot be loaded because running scripts is disabled} & PowerShell blokira skripte &
  Jednom pokrenite \code{Set-ExecutionPolicy -Scope CurrentUser RemoteSigned}\\
\code{uvicorn} nije prepoznat & Virtualno okruženje nije aktivirano &
  Aktivirajte ga; ispred prompta mora pisati \code{(api)} ili \code{(.venv)}. Alternativa
  bez aktivacije: \code{uv run uvicorn main:app --reload}\\
\code{failed to connect to the docker API} ili \code{Cannot connect to the Docker daemon} &
  Docker Desktop nije pokrenut & Pokrenite ga i pričekajte \emph{Engine running}\\
\code{WSL 2 installation is incomplete} ili greška o virtualizaciji &
  Nedostaje WSL~2 ili je virtualizacija isključena &
  U PowerShellu kao administrator pokrenite \code{wsl --install}; ako ne pomogne, uključite
  virtualizaciju VT-x ili AMD-V u BIOS-u\\
\code{Bind for 0.0.0.0:5432 failed: port is already allocated} &
  Port 5432 drži lokalni PostgreSQL ili drugi kontejner &
  \code{docker ps} pokaže kontejner koji treba zaustaviti; inače u \code{docker-compose.yml}
  stavite \code{"5433:5432"}, a u \code{api/main.py} \code{localhost:5433}\\
\code{/health} vraća 500, u terminalu \code{ConnectionRefusedError} &
  Baza nije pokrenuta ili se još diže &
  \code{docker compose ps}; pričekajte nekoliko sekundi i osvježite\\
\code{npm} ili \code{node} nije prepoznat &
  Node.js nije instaliran ili je VS Code otvoren prije instalacije &
  Instalirajte Node.js, pa potpuno zatvorite i ponovno otvorite VS Code; terminal
  u VS Code-u vidi novi PATH tek nakon toga\\
\code{npm install} javlja \code{Unsupported engine} & Prestara verzija Node.js-a &
  Instalirajte Node.js 24 LTS\\
\code{code} nije prepoznat na macOS-u & VS Code nije dodan u PATH &
  U VS Code-u: Command Palette → \emph{Shell Command: Install 'code' command in PATH}\\
\bottomrule
\end{tabularx}
\end{center}

\end{document}
